\section{Historical background}
\label{sec:historical_background}

\noindent This section provides the historical context of our empirical analysis. In Section \ref{sec:emergence_of_gold_clauses}, we trace the emergence of gold clauses to the Civil War and discuss their subsequent widespread adoption in long-term obligations. In Section \ref{sec:abrogation_of_gold_clauses}, we examine the events that led the Roosevelt administration to annul the clauses in 1933. For a comprehensive and masterful account of the abrogation of the gold clauses, see \cite{Edwards2018}.

\subsection{Emergence and adoption of the gold clause}
\label{sec:emergence_of_gold_clauses}

\noindent The gold clause emerged in the United States during the Civil War when gold-backed currency and fiat money (greenbacks) circulated simultaneously. During that period, long-term obligations were issued in both currencies \citep{Payne2025}. According to \cite{Edwards2018}, obligations denominated in gold-backed currency were considered safer since ``the amount to be received in payment at some future date was anchored to the price of gold and, thus, not affected by possible changes in the purchasing power of paper money.''

By the time the US returned to the gold standard in 1879, gold clauses had become a standard feature in most corporate and utility bonds, as well as many farm and house mortgages \citep{Jackson1941, Edwards2018}. However, their adoption was far from universal---\cite{Dam1983} reports that only about fifty-five percent of long-term obligations contained gold clauses.

This selective use of the clause in long-term contracts can be seen among the large industrial firms studied in this paper. While the clause is present in the vast majority of corporate bonds (97\% in our sample), it is entirely absent from preferred shares, another type of long-term obligations commonly used at the time. The exact reason why gold clauses were included in nearly all corporate bonds, present in some mortgages, and absent from preferred shares remains unclear. Despite extensive investigation of contemporary sources and legal studies, we found no definitive explanation for this differential adoption pattern.

Independently of why adoption differed, the record shows that the widespread use of the gold clause in corporate bonds originated from the monetary disruptions of the Civil War. This historical accident would prove consequential 70 years later when the United States abandoned the gold standard in the midst of the Great Depression.

\subsection{Abrogation of the gold clause}
\label{sec:abrogation_of_gold_clauses}

\noindent By the onset of the Great Depression, most major economies were on the gold standard, but the crisis soon forced many to abandon it. The United Kingdom led the way. It ended convertibility on September 21, 1931, after speculative pressures and gold outflows severely undermined confidence in the pound's fixed exchange rate.

In the United States, Britain's exit from the standard prompted widespread concern that the dollar might be next, leading depositors to withdraw funds and hoard gold.\footnote{Both the depositor panic and the Federal Reserve's subsequent discount rate increase are documented in \cite{Friedman1963}.} The Federal Reserve raised its discount rate, which swiftly stemmed gold outflows. While this policy protected gold reserves, the monetary tightening proved disastrous for an already fragile banking system, forcing hundreds of banks to close that month alone.\footnote{According to the Federal Reserve Board, 522 banks suspended operations in October 1931 alone.}

Nonetheless, President Hoover defended these painful measures, describing the gold standard as ``a little short of a sacred formula'' \citep{Eichengreen2000}. Hoover would remain a staunch defender of the standard until the end of his presidency, even as the Depression deepened and deflation reached 10\% in 1932. This commitment to the gold standard would also become an important issue in the 1932 presidential election. 

In a campaign speech in October 1932, Hoover accused his opponent Franklin D. Roosevelt of proposing fiat money as relief, warning that it would bring ``one of the most tragic disasters to... the independence of man'' \citep{Edwards2018}. Roosevelt vigorously denied Hoover's accusations. Pointing to the presence of gold clauses in US Treasury bonds, the future president stated that ``no responsible government would have sold to the country securities payable in gold if it knew that the promise---yes, the covenant---embodied in these securities was as dubious as the President of the United States claims it was.'' Yet, only a few weeks after taking office, Roosevelt would reverse course and ask Congress to void that precise covenant.

This dramatic reversal must be understood in the context of the severe banking crisis that greeted Roosevelt's inauguration on March 4, 1933. The day before, depositors had withdrawn \$250 million in gold and \$150 million in currency from the New York Fed alone, leaving it \$250 million short of required reserves. To stabilize the system and curb hoarding, Roosevelt signed the Emergency Banking Act on March 9, authorizing capital injections that partially reassured markets and restored confidence. However, gold redeposits fell short of the administration's expectations \citep{Edwards2018}.

On April 5, 1933, Roosevelt issued Executive Order 6102, banning private gold ownership and requiring all individuals and businesses to surrender their holdings to the Federal Reserve by May 1 at the official price of \$20.67 per ounce. Treasury Secretary William H. Woodin defended the measure, stating that ``gold in private hoards serves no useful purpose... When added to the stock of the Federal Reserve Banks, it serves as a basis for currency and credit.'' While the order effectively ended domestic convertibility, it left the status of gold exports ambiguous. Roosevelt resolved this ambiguity on April 19 by declaring gold exports illegal, formally taking the United States off the gold standard.

The abandonment of the standard marked a dramatic reversal from positions voiced just weeks earlier. As late as March 6, 1933, Treasury Secretary Woodin had insisted it was ``ridiculous and misleading to say that we are off the gold standard... We are definitely on the gold standard.'' Prior to the suspension of convertibility, exchange rates had remained stable, reflecting confidence in the dollar's gold backing. The suspension shocked markets and triggered immediate reactions in foreign exchange, causing the dollar to quickly depreciate against currencies still tied to gold, such as the French franc. Figure \ref{fig:F2} plots the dollar exchange rate against sterling and the franc \rev{from December 1930 through September 1936}, making this decline apparent. This rapid devaluation created a major domestic problem because of the widespread presence of gold clauses in debt contracts.

At that time, the administration estimated that about \$120 billion of debt---roughly two times GDP, with \$100 billion issued by the private sector---was linked to gold (\cite{Edwards2015}). Dollar depreciation therefore implied a nationwide rise in debt burdens just as many borrowers were already struggling. This issue was well understood at the time. For example, in a speech during the 1932 campaign, President Hoover stated that ``A considerable part of farm mortgages, most of our industrial and all of our Government, most of our State and municipal bonds, and most other long-term obligations are written as payable in gold.'' He added that, if the US had gone off gold, farmers ``would have sold their produce for depreciated currency... but you would have to pay a premium on such of your debts as are written in gold'' \citep{Hoover1932}.

Although the administration initially maintained the official gold price at \$20.67 per ounce, the dollar’s implicit depreciation in international markets prompted the first legal challenges to the enforceability of gold clauses.\footnote{With private holdings forbidden and exports prohibited, the official price of gold only specified the purchase price of newly minted gold.}\kern0.1em\footnote{On May 1, 1933, proceedings began over payment of gold-denominated coupons on private bonds that had been paid in dollars (\cite{Edwards2018}).} To address the issue, the administration asked Congress on May 26 to void the gold clauses in all existing and future contracts. On June 5, Congress passed the Joint Resolution abrogating gold clauses, drawing protests from some members of Congress who argued the abrogation repudiated government promises and violated private contracts.\footnote{Throughout the litigation surrounding the abrogation of the gold clause, payments on bonds containing the clause were not indexed to gold.}

In stark contrast to these legal objections, financial markets rallied in response to the Joint Resolution: \rev{over the nine trading days between May 26 and June 6, on average, exposed firms returned 30.9\%, against 25.4\% for non-exposed firms, a differential of 5.5 percentage points ($t = 3.0$; Table~\ref{tab:event_study}, Panel~A).\footnote{\rev{Size breakpoints are computed over the entire pool of firms. Relative to non-exposed firms, a larger share of exposed firms falls in the small tercile (39\% versus 31\%), and small firms earned the largest returns in this window, so the pooled differential exceeds the average of the tercile differentials.}} The differential is strongest among smaller firms, with small-tercile exposed firms returning 49.0\% against 44.4\% for their non-exposed counterparts, and weaker, though still visible, among the largest (18.3\% against 15.9\%).\footnote{\rev{We hold size fixed through terciles rather than by value-weighting the portfolios: value weights would concentrate the comparison in a handful of firms: the five largest gold-exposed firms alone would carry nearly half of the value-weighted gold portfolio, so a value-weighted differential would reflect a few of the era's giants rather than the average exposed firm.}} The within-tercile differentials, however, are noisier and not statistically significant at conventional levels.} \rev{While the Joint Resolution was a positive development for gold-clause issuers, the outcome} was clouded by the possibility that the abrogation would be ruled unconstitutional, creating uncertainty about whether the implicit debt forgiveness would ultimately be upheld.

As early as May 1933, bondholders filed lawsuits requesting the enforcement of gold clauses. These cases, some concerning private debt contracts and others involving government bonds, moved through the courts during 1933–34. In each, the courts faced the question of whether Congress had violated the Constitution by altering existing contracts.\footnote{\rev{The intermediate legal steps (the first court hearing on the enforceability of gold clauses in May 1933, the lower-court rulings upholding the Joint Resolution in 1934, and the Supreme Court's first grant of certiorari in October 1934) produced no significant differential returns between gold-exposed and non-exposed firms, consistent with these procedural steps being largely anticipated. Internet Appendix Section~\ref{sec:ia_other_events} reports these estimates.}}

\rev{In the fall of 1934, with appeals still pending before the circuit court, the government petitioned the Supreme Court to review \textit{United States v.\ Bankers Trust Co.}\ directly, seeking a prompt and definitive ruling on the gold question rather than further lower-court litigation.\footnote{\rev{A writ of certiorari is the Supreme Court's agreement to review a lower-court case.}} On Monday, November~5, the Court granted certiorari. The following day, the president's party achieved an exceptional midterm result, gaining seats in both chambers of Congress; the landslide was widely read as popular ratification of the New Deal, plausibly raising the perceived durability of the abrogation. The stock market's response over November 5--8 was pronounced: gold-exposed firms returned 7.5\%, against 4.0\% for non-exposed firms, a differential of 3.5 percentage points ($t = 3.3$; Table~\ref{tab:event_study}, Panel~B).\footnote{\rev{Because the exchange was closed on election day, the certiorari grant and the election cannot be separated within this window.}} The differential is concentrated among small firms.}

From January 8–10, 1935, the gold cases were argued before the Supreme Court. Press reports indicated that the government’s legal team performed poorly, struggling to answer many questions from the Justices \citep{Edwards2018}.\footnote{\rev{On January 11, the \textit{Washington Post} reported that the government's counsel had not had smooth going \citep{WP1935}, and the \textit{Chicago Daily Tribune} observed that the government's presentation had suffered palpably in comparison with its adversary's \citep{CDT1935a}. On January 12, the \textit{Chicago Daily Tribune} again pointed to the weakness of the government's argument \citep{CDT1935b}.}} \rev{During the three days of oral arguments (January 8--10), stocks reacted little: exposed firms fell 1.3\%, against 1.4\% for non-exposed firms (Table~\ref{tab:event_study}, Panel~C). News of the government's poor showing reached the public only gradually; many newspapers reported on the argument sessions on January 11 and 12 \citep{Edwards2018}. It is as the news spread that a differential appears: between January 11 and January 17, gold-exposed firms fell 4.5\%, against 3.3\% for non-exposed firms (Table~\ref{tab:event_study}, Panel~D). By January 13, the \textit{New York Times} reported that, amid ``the speculative fever for gold,'' Liberty Bond prices had reached their highest level since 1917 \citep{NYT1935}. Since Liberty Bonds were themselves gold-denominated government obligations, their rising prices reflected investor expectations that the government could lose the case and be forced to honor full gold payments.} The Securities and Exchange Commission even stated it would seek presidential authority to close the exchanges if necessary.

On February 18, 1935, the Supreme Court rendered its decision, upholding the abrogation of gold clauses in private debt. The 5-4 vote reflected the uncertainty that markets had perceived during this period. The four conservative justices issued a scathing dissent: ``The Constitution as many of us understood it... is gone... Because Congress may adopt a system, it does not follow that this may be enforced in violation of existing contracts... Shame and humiliation are upon us now. Moral and financial chaos may be confidently expected.'' 

In stark contrast to the minority's opinion, \rev{the market received the news of the abrogation positively: on the day of the decision, gold-exposed firms returned 3.7\%, against 3.2\% for non-exposed firms (Table~\ref{tab:event_study}, Panel~E).} By the time of the ruling, gold had risen to \$35 per ounce, 69\% above the April 1933 level (Figure \ref{fig:F3}), which explains the market's positive reaction: enforcing gold clauses at that rate would have sharply raised the real debt burden of many firms.\footnote{Inflation had been modest in 1933 (0.8\%) and 1934 (1.5\%), so the 69\% increase in nominal payments would have translated into a roughly 66\% increase in real debt burden.} The 1935 decision was final and eliminated this risk.\footnote{The Court heard four cases: two on private contracts and two on public contracts. For private contracts, it ruled 5–4 that abrogation was constitutional. For public contracts, it ruled 8–1 that abrogation was \textit{un}constitutional, but also ruled that plaintiffs had not suffered recoverable damages, such that no compensation was awarded.}

\rev{Panel F of Table~\ref{tab:event_study} completes the picture: over the full episode, from the eve of the Joint Resolution window through the decision, the cumulative differential between exposed and non-exposed firms is modest. The preceding panels are thus not simply picking up a systematic and persistent return differential between gold-exposed and non-exposed firms, reinforcing our interpretation that the differentials at the specific event dates reflect gold-specific news.}

This short history shows how the legal battle over gold created prolonged uncertainty about firms' future liabilities. For nearly two years, legal challenges to the abrogation accumulated as cases worked their way through the courts. The Supreme Court's 1935 ruling ultimately removed the risk of higher debt burdens by upholding the abrogation. \rev{Moreover, Roosevelt's 1933 decisions to take the country off the gold standard and abrogate the clauses appear to have been largely unexpected, as evidenced by market reactions.} This surprise is consistent with the political rhetoric preceding the crisis: President Hoover's unwavering commitment to the gold standard throughout his presidency and Roosevelt's campaign insistence that he would not violate the gold covenant likely reinforced expectations that these obligations would be honored. Consistent with this interpretation, the next section provides further evidence that firm managers did not systematically adjust the gold content of their liabilities in the years preceding the abrogation.
